% Vietnamese translation for OI Public Library
% Source-PDF: IOI中国国家候选队论文/2007/day2/4.刘雨辰《对拟阵的初步研究》.pdf
\documentclass[11pt]{article}
\usepackage{oipl-vn}

\newcommand{\Oh}{\mathcal{O}}
\newcommand{\dd}{\mathrm{d}}
\newcommand{\rank}{\operatorname{rank}}

\title{Nghiên cứu nhập môn về matroid}
\author{Liu Yuchen\\Trường Trung học số 2 Hàng Châu, Chiết Giang}
\date{2007}

\begin{document}
\maketitle

\origtitle{An introductory study of matroids}
\sourcepath{IOI China candidate team papers / 2007 / day 2 / Liu Yuchen}

\section*{Tóm tắt}

Matroid, còn được gọi theo nghĩa hẹp là một dạng trừu tượng hóa của độc lập
tuyến tính, được Whitney đưa ra vào năm 1935. Matroid là một nội dung quan
trọng trong tối ưu hóa tổ hợp và lý thuyết đồ thị; trong vài thập niên gần
đây nó đã phát triển mạnh thành một lý thuyết rộng và sâu.

Bài viết này giới thiệu sơ bộ về matroid. Phần thứ nhất đưa ra khái niệm hệ
tập con và matroid. Phần thứ hai phát biểu bài toán tối ưu trên matroid và
chứng minh tính đúng đắn của thuật toán tham lam; hai phần này luôn đi kèm
hai ví dụ là ba lô phân số và cây khung nhỏ nhất. Phần thứ ba trình bày một
ứng dụng tối ưu hóa trên matroid trong lập lịch. Phần thứ tư nêu thêm một số
ví dụ, đặc biệt là matroid tuyến tính. Phần mở rộng bàn về trò chơi đóng ngắt
của Shannon; các phụ lục tóm tắt nghịch lý Russell và cách cài đặt nhanh bài
toán lập lịch bằng DSU.

\paragraph{Từ khóa.}
Matroid; thuật toán tham lam; trò chơi đóng ngắt Shannon.

\section*{Dẫn nhập}

Matroid xuất hiện khi Whitney nghiên cứu khái niệm độc lập tuyến tính. Về sau,
người ta nhận ra rằng cùng một kiểu lập luận cũng xuất hiện trong nhiều bài
toán đồ thị và tối ưu hóa tổ hợp. Khi một bài toán có cấu trúc matroid, các
tính chất đã chứng minh cho matroid có thể được dùng trực tiếp cho bài toán
đó, thay vì phải chứng minh lại từ đầu.

Bài viết này đặt trọng tâm vào ba nội dung đầu: định nghĩa, bài toán tối ưu
và một ứng dụng lập lịch. Hai ví dụ xuyên suốt là bài toán ba lô phân số, tức
vật phẩm có thể chia nhỏ, và bài toán cây khung nhỏ nhất. Các phần sau khó
hơn; mục đích chủ yếu là mở rộng tầm nhìn và cho thấy matroid không chỉ là một
mẹo tham lam đơn lẻ.

\section{Hệ tập con và khái niệm matroid}

\paragraph{Định nghĩa.}
Một \term{hệ tập con} là một cặp \(M=(S,L)\) thỏa mãn:
\begin{enumerate}
  \item \(S\) là một tập hữu hạn.
  \item \(L\) là một tập hữu hạn, không rỗng, gồm một số tập con của \(S\).
  \item Tính di truyền: với mọi \(B\in L\) và mọi \(A\subseteq B\), ta có
        \(A\in L\). Do đó \(\varnothing\in L\).
\end{enumerate}
Một hệ tập con được gọi là \term{matroid} nếu nó còn thỏa mãn:
\begin{enumerate}
  \setcounter{enumi}{3}
  \item Tính trao đổi: với mọi \(A,B\in L\) và \(|A|<|B|\), tồn tại
        \(x\in B-A\) sao cho \(A\cup\{x\}\in L\).
\end{enumerate}

Có thể hình dung \(S\) là toàn bộ học sinh trong một lớp và \(L\) là tập các
danh sách nhóm hợp lệ. Tính di truyền nói rằng nếu một nhóm hợp lệ thì mọi
nhóm con của nó cũng hợp lệ. Tính trao đổi nói rằng nếu danh sách \(A\) ngắn
hơn danh sách \(B\), ta luôn có thể lấy một người nằm trong \(B\) nhưng không
nằm trong \(A\) để thêm vào \(A\), và danh sách mới vẫn hợp lệ. Người được thêm
phải đến từ \(B-A\), không được là một phần tử ngoài hai danh sách.

Hai tính chất di truyền và trao đổi là cốt lõi của matroid. Khi mô hình hóa
một bài toán, hai điều kiện đầu thường hiển nhiên do cách chọn \(S\) và \(L\);
phần khó là chứng minh hai tính chất này, đặc biệt là tính trao đổi.

\subsection*{Ví dụ ba lô phân số}

Trước hết biến đổi bài toán: một vật có thể tích \(V\) và giá trị đơn vị \(W\)
được xem như \(V\) vật có thể tích \(1\), mỗi vật có giá trị \(W\). Khi đó mọi
vật đều có thể tích đơn vị. Định nghĩa \(M=(S,L)\) như sau:
\[
  S=\{\text{tất cả các vật đơn vị}\},\qquad
  L=\{X\subseteq S: |X|\le \mathrm{MAX}\},
\]
trong đó \(\mathrm{MAX}\) là dung tích ba lô.

Hệ này có tính di truyền vì \(Y\subseteq X\) và \(|X|\le\mathrm{MAX}\) kéo theo
\(|Y|\le\mathrm{MAX}\). Nếu \(A,B\in L\) và \(|A|<|B|\), chọn tùy ý
\(x\in B-A\). Khi đó \(|A\cup\{x\}|=|A|+1\le |B|\le \mathrm{MAX}\), nên
\(A\cup\{x\}\in L\). Vì vậy đây là một matroid; ta gọi nó là matroid của bài
toán ba lô.

\subsection*{Ví dụ đồ thị và cây khung nhỏ nhất}

Với một đồ thị vô hướng \(G=(V,E)\), định nghĩa:
\[
  S=E,\qquad L=\{X\subseteq E: \text{đồ thị do } X \text{ tạo ra không có chu trình}\}.
\]
Tính di truyền là hiển nhiên: xóa bớt cạnh không thể tạo thêm chu trình.

Ta chứng minh tính trao đổi. Lấy \(A,B\in L\) với \(|A|<|B|\). Gọi \(G_A\) và
\(G_B\) là hai rừng do \(A\) và \(B\) tạo ra. Một rừng trên \(|V|\) đỉnh và
\(k\) cạnh có \(|V|-k\) thành phần liên thông. Do \(|A|<|B|\), rừng \(G_B\) có
ít thành phần liên thông hơn \(G_A\). Vì vậy tồn tại một thành phần liên thông
\(T\) của \(G_B\) mà các đỉnh của nó không cùng nằm trong một thành phần liên
thông của \(G_A\). Trong \(T\) có một cạnh \(x\in B\) nối hai thành phần liên
thông khác nhau của \(G_A\). Khi thêm \(x\) vào \(A\), ta không tạo chu trình.
Hơn nữa \(x\notin A\), nên \(x\in B-A\) và \(A\cup\{x\}\in L\). Do đó \(M\) là
một matroid, gọi là matroid đồ thị của \(G\).

\subsection*{Một số thuật ngữ}

Với \(U\subseteq S\), nếu \(U\in L\) thì \(U\) được gọi là một \term{tập độc lập}.
Một tập độc lập \(A\) được gọi là \term{mở rộng được} nếu tồn tại
\(x\in S-A\) sao cho \(A\cup\{x\}\in L\). Tập độc lập không mở rộng được gọi là
\term{tập độc lập cực đại}. Với \(U\subseteq S\), định nghĩa
\[
  r(U)=\max\{|X|: X\subseteq U,\ X\in L\}.
\]
Giá trị \(r(U)\) gọi là \term{hạng} của \(U\), và \(r\) là hàm hạng của
matroid.

\paragraph{Định lý 1.}
Mọi tập độc lập cực đại của một matroid đều có cùng kích thước.

\paragraph{Chứng minh.}
Giả sử \(A\) và \(B\) là hai tập độc lập cực đại. Nếu \(|A|\ne |B|\), không mất
tính tổng quát giả sử \(|A|<|B|\). Theo tính trao đổi, tồn tại \(x\in B-A\) sao
cho \(A\cup\{x\}\in L\). Khi đó \(A\) mở rộng được, mâu thuẫn với tính cực đại.

\section{Bài toán tối ưu trên matroid}

Với một matroid \(M=(S,L)\), gán cho mỗi phần tử \(x\in S\) một trọng số nguyên
dương \(w(x)\). Trọng số của tập con \(U\subseteq S\) là
\[
  w(U)=\sum_{x\in U}w(x).
\]
Bài toán tối ưu trên matroid yêu cầu tìm một tập độc lập có trọng số lớn nhất.
Vì trọng số đều dương, một nghiệm tối ưu luôn là tập độc lập cực đại.

Ở ví dụ ba lô, \(w(x)\) là giá trị của vật đơn vị \(x\). Bài toán trở thành tìm
tập vật đơn vị độc lập có tổng giá trị lớn nhất, tức một phương án lấy đồ có
giá trị tối đa.

Ở ví dụ cây khung nhỏ nhất, mục tiêu ban đầu là tối thiểu hóa tổng trọng số
cạnh. Ta chuyển nó thành tối đa hóa bằng cách đặt
\[
  w(x)=-g(x)+\Delta,\qquad \Delta=\max_{e\in E}g(e)+1,
\]
trong đó \(g(x)\) là trọng số cạnh gốc. Khi đó \(w(x)>0\). Vì mọi cây khung có
cùng số cạnh, việc tối đa hóa tổng \(w\) tương đương với tối thiểu hóa tổng
\(g\).

\subsection*{Thuật toán tham lam}

Thuật toán tham lam cho matroid có trọng số:
\begin{verbatim}
Greedy(M, w)
    A := empty set
    sort S[M] in nonincreasing order of w
    for each x in S[M] in that order do
        if A union {x} belongs to L[M] then
            A := A union {x}
    return A
\end{verbatim}

Với ba lô phân số, thuật toán chọn các vật đơn vị theo giá trị giảm dần và bỏ
vào ba lô nếu còn chỗ. Với cây khung nhỏ nhất, vì trọng số đã đổi dấu và tịnh
tiến, thuật toán xét các cạnh theo trọng số gốc tăng dần; nếu thêm cạnh không
tạo chu trình thì nhận cạnh đó. Đây chính là tinh thần của thuật toán Kruskal.

Nếu \(n=|S|\), bước sắp xếp tốn \(\Theta(n\log n)\). Bước tham lam thực hiện
\(\Theta(n)\) lần kiểm tra độc lập. Nếu một lần kiểm tra tốn
\(\Theta(f(n))\), tổng thời gian là
\[
  \Theta(n\log n+n f(n)).
\]
Nút thắt nằm ở cách kiểm tra tập độc lập. Các bài toán khác nhau có cùng cấu
trúc matroid, nhưng phần kiểm tra độc lập thường mang đặc trưng riêng của từng
bài.

\subsection*{Tính đúng đắn}

Ta chứng minh rằng sau mỗi bước, tập \(A\) do thuật toán xây dựng luôn là tập
con của một nghiệm tối ưu nào đó. Ban đầu \(A=\varnothing\), điều này hiển
nhiên đúng.

\paragraph{Định lý 2.}
Cho matroid \(M=(S,L)\). Giả sử \(A\in L\), \(A\subseteq T\), trong đó \(T\) là
một tập độc lập tối ưu. Giả sử \(x\) thỏa mãn \(A\cup\{x\}\in L\), và trong số
các phần tử \(y\) có thể mở rộng \(A\), không có phần tử nào có
\(w(y)>w(x)\). Đặt \(A'=A\cup\{x\}\). Khi đó tồn tại một nghiệm tối ưu chứa
\(A'\).

\paragraph{Chứng minh.}
Ta chứng minh bằng phản chứng và trao đổi. Nếu \(x\in T\) thì \(A'\subseteq T\)
và không có gì phải chứng minh. Giả sử \(x\notin T\). Bắt đầu với
\(T'=A'\). Khi \(|T'|<|T|\), áp dụng tính trao đổi cho \(T'\) và \(T\), ta có
thể chọn một phần tử \(z\in T-T'\) sao cho \(T'\cup\{z\}\in L\), rồi thêm
\(z\) vào \(T'\). Lặp lại đến khi \(|T'|=|T|\).

Khi kết thúc, \(T'\) là một tập độc lập chứa \(A'\) và có cùng kích thước với
\(T\). Nó có dạng
\[
  T'=(T-\{y\})\cup\{x\}
\]
với một phần tử \(y\in T-A\) không được thêm lại. Vì \(y\) cũng là một phần tử
có thể mở rộng \(A\), cách chọn \(x\) cho ta \(w(y)\le w(x)\). Do đó
\[
  w(T')=w(T)-w(y)+w(x)\ge w(T).
\]
Nếu bất đẳng thức nghiêm, \(T\) không tối ưu, mâu thuẫn. Nếu bằng nhau, \(T'\)
cũng là nghiệm tối ưu và chứa \(A'\), trái với giả thiết rằng không nghiệm tối
ưu nào chứa \(A'\). Vậy luôn tồn tại nghiệm tối ưu chứa \(A'\).

Chứng minh trên làm rõ vai trò của tính trao đổi. Nếu chỉ có một hệ tập con
có tính di truyền thì chưa đủ để tham lam đúng; điều kiện trao đổi chính là
phần giúp ta thay thế phần tử trong nghiệm tối ưu mà không làm giảm chất
lượng nghiệm. Chứng minh quen thuộc của Kruskal cũng là một trường hợp cụ thể
của lập luận này: thêm cạnh đang xét vào một cây khung tối ưu sẽ tạo chu trình,
rồi bỏ khỏi chu trình một cạnh không tốt hơn cạnh vừa thêm.

\section{Một bài toán lập lịch}

Xét bài toán lập lịch một bộ xử lý cho các công việc thời gian đơn vị. Có tập
công việc \(S=\{1,2,\ldots,n\}\). Mỗi công việc \(i\) có hạn chót
\(d_i\), \(1\le d_i\le n\), và tiền phạt dương \(w_i\). Một lịch là một hoán
vị của \(S\): công việc đầu chạy trong khoảng \([0,1]\), công việc thứ hai
trong \([1,2]\), và cứ thế. Nếu công việc \(i\) kết thúc sau thời điểm \(d_i\),
ta phải trả phạt \(w_i\). Mục tiêu là tối thiểu hóa tổng phạt.

Với một lịch bất kỳ, ta chỉ cần quan tâm công việc nào hoàn thành đúng hạn.
Các công việc trễ có thể đưa về cuối lịch mà không làm xấu hơn nghiệm. Vì vậy
ta xét câu hỏi: với một tập \(A\subseteq S\), có tồn tại lịch sao cho mọi công
việc trong \(A\) đều hoàn thành đúng hạn không? Nếu có, gọi \(A\) là
\term{khả thi}.

Với một tập khả thi \(A\) và số nguyên \(t\), đặt \(F(A,t)\) là số công việc
trong \(A\) có hạn chót không vượt quá \(t\).

\paragraph{Định lý 3.}
Ba mệnh đề sau tương đương:
\begin{enumerate}
  \item \(A\) là khả thi.
  \item Với mọi \(t=1,2,\ldots,n\), ta có \(F(A,t)\le t\).
  \item Lịch sắp các công việc trong \(A\) theo hạn chót tăng dần hoàn thành
        đúng hạn toàn bộ \(A\).
\end{enumerate}

\paragraph{Chứng minh.}
Nếu \(A\) khả thi mà tồn tại \(t\) với \(F(A,t)>t\), thì hơn \(t\) công việc có
hạn chót không quá \(t\) phải được hoàn thành trong \(t\) đơn vị thời gian đầu,
điều không thể. Do đó mệnh đề 1 kéo theo mệnh đề 2.

Nếu mệnh đề 2 đúng, xét lịch theo hạn chót tăng dần. Một công việc có hạn chót
\(t\) hoàn thành không muộn hơn thời điểm \(F(A,t)\), mà \(F(A,t)\le t\). Vậy
mệnh đề 2 kéo theo mệnh đề 3. Mệnh đề 3 hiển nhiên kéo theo mệnh đề 1.

Định nghĩa \(M=(S,L)\), trong đó
\[
  L=\{X\subseteq S: X \text{ là khả thi}\}.
\]
Tính di truyền rõ ràng đúng: nếu một tập công việc làm đúng hạn được thì mọi
tập con của nó cũng làm đúng hạn được.

Ta chứng minh tính trao đổi. Lấy hai tập khả thi \(A,B\) với \(|A|<|B|\). Chọn
\(k\) lớn nhất sao cho \(F(B,k)\le F(A,k)\). Vì
\(F(B,n)=|B|>|A|=F(A,n)\), nên \(k<n\); vì \(F(B,0)=F(A,0)=0\), nên \(k\ge 0\).
Từ \(F(B,k)\le F(A,k)\) và \(|B|>|A|\), số công việc trong \(B\) có hạn chót
lớn hơn \(k\) nhiều hơn số công việc như vậy trong \(A\). Vì thế tồn tại
\(x\in B-A\) có hạn chót \(t>k\). Đặt \(A'=A\cup\{x\}\).

Với \(s<t\), ta có \(F(A',s)=F(A,s)\le s\). Với \(s\ge t\), do cách chọn \(k\),
ta có \(F(B,s)>F(A,s)\), nên
\[
  F(A',s)=F(A,s)+1\le F(B,s)\le s.
\]
Theo Định lý 3, \(A'\) khả thi. Vậy \(M\) là một matroid.

Bài toán tối thiểu hóa phạt tương đương với tối đa hóa tổng tiền phạt của các
công việc được làm đúng hạn. Do \(M\) là matroid, ta sắp công việc theo
\(w_i\) giảm dần và lần lượt nhận một công việc nếu tập đã nhận vẫn khả thi.
Cài đặt thô kiểm tra điều kiện \(F(A,s)\le s\) trong \(\Theta(n)\), cho tổng
thời gian \(\Theta(n^2)\) sau khi sắp xếp. Có thể cài đặt nhanh hơn bằng DSU,
được tóm tắt ở Phụ lục II.

\section{Một số ví dụ về matroid}

\subsection*{Matroid tuyến tính}

Cho một ma trận \(T\) kích thước \(m\times n\) trên một trường, chẳng hạn trên
\(\mathbb{R}\). Đặt \(S\) là tập các cột của \(T\). Một tập con \(U\subseteq S\)
được gọi là độc lập khi các cột trong \(U\) độc lập tuyến tính. Khi đó
\(M=(S,L)\) là một matroid, gọi là \term{matroid tuyến tính}.

Tính di truyền là tính chất cơ bản của độc lập tuyến tính: mọi tập con của một
hệ vector độc lập tuyến tính vẫn độc lập tuyến tính. Với tính trao đổi, lấy
hai tập độc lập \(X,Y\) và \(|Y|<|X|\). Gọi \(W\) là không gian vector sinh bởi
\(X\cup Y\). Khi đó \(\dim W\ge |X|\). Nếu với mọi \(x\in X-Y\), tập
\(Y\cup\{x\}\) đều phụ thuộc tuyến tính, thì mọi vector trong \(X-Y\) đều thuộc
không gian sinh bởi \(Y\). Suy ra \(W\) cũng do \(Y\) sinh ra, nên
\(\dim W\le |Y|\), mâu thuẫn với \(|X|>|Y|\). Do đó tồn tại \(x\in X-Y\) sao
cho \(Y\cup\{x\}\) độc lập.

Một bài toán tối ưu trên matroid tuyến tính cần kiểm tra liệu một hệ vector có
độc lập tuyến tính hay không. Việc này tương đương với tính hạng ma trận, và
thường được cài đặt bằng khử Gauss thay vì định thức trực tiếp.

Matroid đồ thị cũng có thể xem là một matroid tuyến tính. Với đồ thị vô hướng
\(G=(V,E)\), định nghĩa ma trận đỉnh-cạnh \(B\) kích thước \(|V|\times |E|\).
Mỗi cột ứng với một cạnh, có đúng hai phần tử khác \(0\), một bằng \(1\) và
một bằng \(-1\), ở hai hàng là hai đầu mút của cạnh. Một tập cạnh không tạo
chu trình khi và chỉ khi các cột tương ứng độc lập tuyến tính. Vì vậy matroid
đồ thị là một trường hợp của matroid tuyến tính. Chi tiết chứng minh khá dài
nên bản gốc lược bỏ.

\subsection*{Matroid đều}

Với một tập hữu hạn \(S\), đặt \(I_k\) là tập tất cả tập con của \(S\) có kích
thước không quá \(k\). Khi đó \(M=(S,I_k)\) là một matroid, thường gọi là
matroid đều. Matroid của ví dụ ba lô sau khi tách vật phẩm về thể tích đơn vị
chính là một matroid đều.

\subsection*{Matroid đối ngẫu}

Nếu \(M=(S,L)\) là một matroid, định nghĩa \(M'=(S,L')\) bởi
\[
  L'=\{A'\subseteq S:\text{ tồn tại một tập độc lập cực đại } A
      \text{ của } M \text{ sao cho } A\subseteq S-A'\}.
\]
Ta gọi \(M'\) là \term{matroid đối ngẫu} của \(M\).

Tính di truyền: nếu \(B'\in L'\) và \(A'\subseteq B'\), tồn tại một tập độc lập
cực đại \(X\) của \(M\) sao cho \(X\subseteq S-B'\). Khi đó
\(X\subseteq S-A'\), nên \(A'\in L'\).

Với tính trao đổi, lấy \(A',B'\in L'\) và \(|A'|<|B'|\). Tồn tại các tập độc
lập cực đại \(A,B\) của \(M\) sao cho \(A\subseteq S-A'\) và
\(B\subseteq S-B'\). Chọn \(x\in B'-A'\), đặt \(C'=A'\cup\{x\}\). Nếu
\(x\notin A\) thì \(A\subseteq S-C'\) và \(C'\in L'\). Trường hợp khó là
\(x\in A\). Khi đó dùng tính trao đổi trong \(M\) để thay \(x\) trong \(A\) bởi
một phần tử \(y\in B-A\), thu được một tập độc lập cực đại
\(C=A-\{x\}\cup\{y\}\) vẫn nằm trong \(S-C'\). Vậy \(C'\in L'\), và \(M'\) là
một matroid.

\subsection*{Matroid ghép cặp}

Với đồ thị vô hướng \(G=(V,E)\), đặt \(S=V\). Một tập đỉnh \(A\subseteq V\)
được gọi là độc lập nếu tồn tại một matching của \(G\) phủ toàn bộ các đỉnh
trong \(A\). Nếu \(G\) có matching hoàn hảo thì mọi tập con của \(V\) đều độc
lập, nhưng đó chỉ là một trường hợp đặc biệt.

Hệ này là một matroid. Tính di truyền rõ ràng: một matching phủ \(A\) cũng phủ
mọi tập con của \(A\). Tính trao đổi là phần khó. Lấy hai tập độc lập \(A,B\)
với \(|A|<|B|\), và hai matching \(P_A,P_B\) lần lượt phủ \(A,B\). Nếu \(P_A\)
đã phủ một đỉnh trong \(B-A\), ta thêm ngay đỉnh đó vào \(A\). Nếu không, xét
hiệu đối xứng \(P_A\triangle P_B\). Từ mỗi đỉnh trong \(B-A\) có một đường đi
luân phiên theo cạnh của \(P_A\) và \(P_B\). Không phải mọi đường như vậy đều
có thể kết thúc trong \(A-B\), vì \(|B-A|>|A-B|\). Do đó tồn tại một đường đi
luân phiên bắt đầu ở một đỉnh của \(B-A\) và kết thúc ngoài \(A\). Đảo trạng
thái các cạnh trên đường này tạo một matching mới phủ \(A\) và thêm một đỉnh
thuộc \(B-A\). Đây chính là tính trao đổi. Bản gốc ghi chú rằng chứng minh này
khá cô đọng và thực chất là lập luận đường tăng.

\section{Mở rộng: trò chơi đóng ngắt Shannon}

Trò chơi đóng ngắt Shannon do Claude Shannon đề xuất; lời giải đẹp được
A. Lehman tìm ra. Bản gốc chỉ trích lược một phần vì các tài liệu mà tác giả
tìm được khá rời rạc và chứng minh đầy đủ của một định lý nền tảng rất phức
tạp.

Trò chơi diễn ra trên một đồ thị vô hướng \(G=(V,E)\), có thể có cạnh song
song. Có hai người chơi: người nối và người cắt. Mỗi lượt, người nối chọn một
cạnh chưa được đánh dấu và đánh dấu \(+\); người cắt chọn một cạnh chưa được
đánh dấu \(+\) và xóa nó. Người cắt đi trước, hai bên luân phiên cho đến khi
không thể đi tiếp. Khi kết thúc, nếu đồ thị còn liên thông thì người nối
thắng; nếu không thì người cắt thắng.

Điều kiện cần và đủ để người nối, dù đi sau, vẫn thắng là: đồ thị chứa hai
cây khung độc lập cạnh, tức hai cây khung không có cạnh chung. Nếu xét matroid
đồ thị \(M=(S,L)\) với \(S=E\), điều kiện này có thể viết là tồn tại hai tập
độc lập cực đại \(A,B\) sao cho \(A\cap B=\varnothing\).

\subsection*{Một sự kiện về hợp matroid}

Cho \(k\) matroid \(M_i=(S_i,L_i)\). Hợp của chúng là \(M=(S,L)\), trong đó
\[
  S=\bigcup_{i=1}^k S_i,\qquad
  L=\{I_1\cup I_2\cup\cdots\cup I_k: I_i\in L_i\}.
\]
Hợp này cũng là một matroid. Chứng minh khá phức tạp nên bản gốc không trình
bày. Nếu \(r_i\) là hàm hạng của \(M_i\), hàm hạng \(r\) của \(M\) thỏa mãn
\[
  r(U)=\min_{T\subseteq U}\left(|U-T|+\sum_{i=1}^k r_i(T\cap S_i)\right).
\]

Áp dụng cho \(k\) bản sao của cùng một matroid \(M=(S,L)\), gọi \(R\) là hàm
hạng của hợp, còn \(r\) là hàm hạng gốc. Khi đó
\[
  R(U)=\min_{T\subseteq U}\bigl(|U-T|+k r(T)\bigr).
\]
Vì thế \(M\) có \(k\) tập độc lập đôi một rời nhau phủ đủ hạng khi và chỉ khi,
với mọi \(T\subseteq S\),
\[
  |S-T|\ge k\bigl(r(S)-r(T)\bigr).
\]
Kết luận này được dùng trong chứng minh dưới đây.

\subsection*{Điều kiện thắng}

\paragraph{Tính cần thiết.}
Nếu không tồn tại hai tập độc lập rời nhau, thì theo kết luận trên với \(k=2\),
tồn tại \(T\subseteq S\) sao cho
\[
  |S-T|<2\bigl(r(S)-r(T)\bigr).
\]
Chiến lược của người cắt là luôn xóa một cạnh trong \(S-T\). Người cắt có thể
xóa ít nhất \(\lceil |S-T|/2\rceil\) cạnh ở phần này, còn người nối cứu được
nhiều nhất \(\lfloor |S-T|/2\rfloor\) cạnh ở phần này. Do bất đẳng thức trên,
số cạnh hữu dụng mà người nối giữ được trong \(S-T\) nhỏ hơn \(r(S)-r(T)\);
trong \(T\), người nối giữ được nhiều nhất \(r(T)\) cạnh hữu dụng. Tổng số
cạnh hữu dụng được giữ nhỏ hơn \(r(S)=|V|-1\), nên cuối cùng đồ thị không liên
thông. Vì vậy người nối thua.

\paragraph{Tính đủ.}
Giả sử đồ thị có hai cây khung rời cạnh \(T_1,T_2\). Ý tưởng là sau mỗi lượt
của người cắt và lượt đáp lại của người nối, ta xây dựng lại một cặp cây
khung mới có nhiều hơn trước một cạnh chung, và mọi cạnh chung đó đều đã được
đánh dấu \(+\).

Ban đầu đặt \(T_{10}=T_1\), \(T_{20}=T_2\). Ở lượt đầu, người cắt xóa một cạnh
\(\beta\).

Nếu \(\beta\) thuộc một trong hai cây, giả sử thuộc \(T_{10}\), thì vì
\(T_{20}\) là cây khung, tồn tại một cạnh \(\alpha\in T_{20}\) sao cho thêm
\(\alpha\) vào \(T_{10}\) rồi bỏ \(\beta\) vẫn thu được một cây khung \(T_{11}\).
Người nối đánh dấu \(+\) trên \(\alpha\), và đặt \(T_{21}=T_{20}\). Khi đó hai
cây mới có đúng một cạnh chung, là \(\alpha\), và cạnh đó đã được đánh dấu.

Nếu \(\beta\) không thuộc cây nào, người nối chọn một cạnh \(\alpha\) bất kỳ
trong \(T_{10}\cup T_{20}\), chẳng hạn \(\alpha\in T_{10}\), rồi đánh dấu nó.
Vì \(T_{20}\) là cây khung, có một cạnh \(\gamma\in T_{20}\) sao cho thêm
\(\alpha\) vào \(T_{20}\) rồi bỏ \(\gamma\) vẫn là cây khung. Đặt
\(T_{11}=T_{10}\) và gọi cây mới là \(T_{21}\). Hai cây mới lại có đúng một
cạnh chung đã đánh dấu.

Các lượt sau lặp lại cùng ý tưởng. Sau lượt thứ \(k\), ta duy trì hai cây
khung \(T_{1k},T_{2k}\) có đúng \(k\) cạnh chung và mọi cạnh chung đều mang
dấu \(+\). Nếu đồ thị có \(m\) đỉnh, sau \(m-1\) lượt hai cây có \(m-1\) cạnh
chung, tức chúng là cùng một cây khung. Các cạnh mang dấu \(+\) vì thế tạo
thành một cây khung, nên đồ thị còn liên thông và người nối thắng.

\section*{Tổng kết}

Matroid là một cấu trúc tổ hợp có hai tính chất bản chất: di truyền và trao
đổi. Khi chứng minh một cấu trúc là matroid, di truyền thường dễ, còn trao đổi
mới là điểm mấu chốt. Trong các chứng minh ở trên, hai kỹ thuật xuất hiện liên
tục là phản chứng và xây dựng đối tượng thay thế.

Tính đúng đắn của thuật toán tham lam trên matroid đến từ chính tính trao đổi.
Matroid ban đầu được phát hiện từ độc lập tuyến tính, và matroid đồ thị cũng
có thể nhìn như một matroid tuyến tính. Trò chơi đóng ngắt Shannon cho thấy
matroid còn liên hệ với các bài toán trò chơi và liên thông; một số nội dung
bị lược chứng minh trong phần mở rộng liên quan đến các thuật toán sâu hơn
trên matroid.

\section*{Tài liệu tham khảo}

\begin{enumerate}
  \item Thomas H. Cormen, Charles E. Leiserson, Ronald L. Rivest, Clifford
        Stein, \textit{Introduction to Algorithms}.
  \item MIT, \textit{Topics in Combinatorial Optimization}.
  \item Richard A. Brualdi, \textit{Introductory Combinatorics}.
  \item Liu Rujia, Huang Liang, \textit{The Art of Algorithms and Informatics
        Competitions}.
  \item James Oxley, \textit{What is a matroid?}
  \item Richard Courant, Herbert Robbins, \textit{What Is Mathematics?}
  \item Lai Hongjian, \textit{Matroid Theory}.
  \item Course materials on linear algebra, Tongji University.
  \item Richard Mansfield, \textit{Strategies for the Shannon switching game}.
  \item Michael Artin, \textit{Algebra}.
\end{enumerate}

\appendix

\section{Nghịch lý Russell}

Phụ lục gốc dùng nghịch lý Russell để nhắc rằng không thể dùng khái niệm
``tập hợp'' một cách hoàn toàn không ràng buộc. Phần tóm tắt như sau.

Gọi một tập là ``thông thường'' nếu nó không chứa chính nó làm phần tử. Chẳng
hạn tập các số nguyên chỉ chứa số nguyên, nên nó không chứa chính nó. Một số
mô tả tập hợp tự quy có thể tạo ra tập chứa chính nó, và ta gọi chúng là
``không thông thường''. Bây giờ xét \(C\), tập của tất cả các tập thông
thường. Nếu \(C\) thông thường, theo định nghĩa của \(C\), nó phải chứa chính
nó, vậy nó lại không thông thường. Nếu \(C\) không thông thường, tức \(C\) chứa
chính nó, thì \(C\) chứa một tập không thông thường, trái với định nghĩa rằng
nó chỉ chứa các tập thông thường. Hai khả năng đều mâu thuẫn. Vì vậy, chỉ
riêng việc giả sử \(C\) tồn tại đã buộc ta phải công nhận các hạn chế và tiên
đề chặt chẽ hơn trong lý thuyết tập hợp.

\section{Cài đặt nhanh bài toán lập lịch}

Thuật toán tham lam xét công việc theo tiền phạt giảm dần. Với công việc
\(j\), nếu còn một ô thời gian trống kết thúc không muộn hơn \(d_j\), ta đặt
\(j\) vào ô trống muộn nhất như vậy; nếu không còn ô trống, công việc này bị
phạt.

Ta duy trì các thời điểm bằng DSU. Hai thời điểm \(i,j\) thuộc cùng một lớp
tương đương khi mọi ô thời gian giữa chúng đã bị lấp đầy. Đại diện của mỗi
lớp được giữ là thời điểm nhỏ nhất trong lớp. Khi lấp ô \((k-1,k]\), ta hợp
hai thời điểm \(k-1\) và \(k\). Với một công việc có hạn chót \(d_j\), nếu
\(\mathrm{find}(d_j)=0\) thì mọi ô trước hạn chót đã đầy; ngược lại, ô trống
muộn nhất là ô kết thúc tại \(\mathrm{find}(d_j)\).

\begin{verbatim}
fillchar(f, sizeof(f), 255)
ret := 0
for i := 1 to n do
begin
    k := find(d[i])
    if k > 0 then
        Union(k - 1, k)
    else
        ret := ret + w[i]
end

procedure Union(u, v: longint);
begin
    u := find(u); v := find(v)
    if u > v then f[u] := v
             else f[v] := u
end
\end{verbatim}

\end{document}
